\section{From potential comparisons to a labelling problem}
\label{sec:boxes}

This section specifies the finite-box input used from the companion,
then proves the game reduction at the width required by our recursion.
The outcome is a labelling problem whose obligations are expressed by
potential vectors. Later sections solve that problem without finding
the potential that certifies the winning partition.

\subsection{The common finite-box interface}

Fix the input graph and ownership throughout. Following the integer
cycle perturbation of Bj\"orklund and Vorobyov
\cite[Section~9]{BjorklundVorobyov2007}, put
\begin{equation}
 \widehat w(e)=(n+1)w(e)+1,
 \qquad W=\max_{e\in E}|\widehat w(e)|.
 \label{eq:bias}
\end{equation}
The edge set is nonempty. Each transformed weight is a nonzero integer,
so $W\ge1$. If a simple directed cycle has length $\ell$ and
original sum $c$, then
\begin{equation}
 \widehat w(C)=(n+1)c+\ell,\qquad 1\le\ell\le n.
 \label{eq:cycle-signs}
\end{equation}
For $c\ge0$ the transformed sum is positive; for $c\le-1$ it is
negative. Thus no simple cycle has zero transformed sum, and positive
transformed sums correspond exactly to nonnegative original sums.
A self-loop has $\ell=1$, and the calculation applies to every
actual edge sequence when parallel edges are present.

For all $z\in\Z^V$, define
\begin{equation}
 F_i(z)=
 \begin{cases}
  \displaystyle\max_{e:i\to j}(\widehat w(e)+z_j),&i\in V_{\mathrm{Max}},\\[2mm]
  \displaystyle\min_{e:i\to j}(\widehat w(e)+z_j),&i\in V_{\mathrm{Min}}.
 \end{cases}
 \label{eq:operator}
\end{equation}
The finite, nonempty outgoing-edge sets ensure that these extrema are
attained. With coordinatewise vector order and common scalar shifts,
\begin{equation}
 z\le z'\ \Longrightarrow\ F(z)\le F(z'),
 \qquad F(z+t)=F(z)+t\quad(t\in\Z).
 \label{eq:operator-properties}
\end{equation}
For integer vectors $A\le u$, the notation $[A,u]$ denotes all
integer vectors between them. A \emph{subsolution} is a full vector
$x\in[A,u]$ satisfying $x_i\le F_i(x)$ at every coordinate where
$x_i>A_i$. A \emph{supersolution} is a full vector $y\in[A,u]$
satisfying $y_i\ge F_i(y)$ wherever $y_i<u_i$. No inequality is
required on the respective excluded boundary. These conventions agree
with the companion's Definition~2.2~\cite{OpenAI2026Randomized}.

\begin{lemma}[Finite-box comparison and basic iteration]
\label{lem:accepted-box}
For the operator in Equation~\eqref{eq:operator} and arbitrary finite
integer bounds $A\le u$, every subsolution $x$ and supersolution
$y$ in $[A,u]$ satisfy $x_i\le y_i$ for all $i\in V$. There is a
unique $h\in[A,u]$ that is both a subsolution and a supersolution.
It is precisely the unique fixed point of
\[
 T_{A,u}(z)=\min\{u,\max\{A,F(z)\}\}.
\]
Consequently every such pair satisfies $x\le h\le y$.

The following synchronous iteration finds $h$. Start with $p=A$;
evaluate $q=T_{A,u}(p)$; return $p$ if $q=p$, and otherwise set
$p=q$ and repeat. It uses at most
\[
                         1+\sum_{i\in V}(u_i-A_i)
\]
evaluations of $T_{A,u}$, including the final stability test.
\end{lemma}

We use this statement from the complete September 25 companion,
Lemmas~2.3--2.4~\cite{OpenAI2026Randomized}. It permits negative
endpoints and collapsed coordinates $A_i=u_i$. In particular its
comparison conclusion concerns all vertices, even though each
hypothesis is imposed only on its stated support. The operator and
integer hypotheses are exactly those above, and
Equation~\eqref{eq:cycle-signs} verifies the required absence of zero
simple-cycle sums. The sandwich inequality follows by comparing with
$h$ on each side. We call $h$ the \emph{box solution} and the
iteration $\Basic$. The algorithm will invoke $\Basic$ only when
the box width is bounded by an absolute constant.

\subsection{The two bands of a wide box}

Choose the initial dyadic width
\begin{equation}
 D_0=\min\{2^d:d\in\Z_{\ge0},\ 2^d\ge64(n+1)(W+1)\}.
 \label{eq:initial-width}
\end{equation}
The next proposition establishes the game reduction for this width.
The companion uses width $8nW$ in its own reduction; that different
choice is not an input to the proof here.

\begin{proposition}[Winning partition from the wide box]
\label{prop:game-reduction}
Let $h$ be the box solution in $[0,D_0]^V$. Then
\[
 V_{\lo}=\{i:h_i\le nW\},\qquad
 V_{\hi}=\{i:h_i\ge D_0-nW\}
\]
partition $V$. There is a positional Max strategy under which every
play starting in $V_{\hi}$ has original mean payoff at least zero.
There is a positional Min strategy under which every play starting
in $V_{\lo}$ has original upper limiting average at most $-1/n$.
Both assertions allow arbitrary history-dependent opposition.
\end{proposition}

\begin{proof}
We first locate the coordinates of $h$, then construct the two
strategies. If $0<h_i<D_0$, both box inequalities apply and give
$h_i=F_i(h)$. Follow an edge attaining that equality and continue
from its head as long as the current coordinate is interior. A first
repeated vertex would yield a simple cycle of tight equalities
$h_i=\widehat w(e)+h_j$, with transformed sum zero. This is
impossible by Equation~\eqref{eq:cycle-signs}. A boundary is
therefore reached within $n$ edges. Each edge changes the potential
by at most $W$. Including the length-zero path from a boundary
coordinate, every coordinate is within $nW$ of either $0$ or
$D_0$. Since $D_0>2nW$, the two bands are disjoint and exhaust $V$.

Every high coordinate is above the lower boundary. At a high Max
vertex choose one outgoing edge satisfying
\begin{equation}
                         h_i\le\widehat w(e)+h_j;
 \label{eq:high-edge}
\end{equation}
such an edge exists by the subsolution inequality. At a high Min
vertex every outgoing edge satisfies this inequality. Each of these
permitted edges has
\[
                  h_j\ge D_0-(n+1)W>nW,
\]
so its head is high by the partition just proved. Fixing the Max
choices and extending them arbitrarily outside the high band gives
a positional strategy keeping all consistent plays from high
vertices in that band.

On a simple cycle of permitted edges,
Equation~\eqref{eq:high-edge} telescopes to a nonnegative transformed
sum. Its original sum is therefore nonnegative. Let
$W_0=\max_e|w(e)|$. Delete simple cycles successively from any
finite play prefix until fewer than $n$ edges remain. The removed
cycles have nonnegative original sums, and the remaining walk has
weight at least $-(n-1)W_0$. Hence every prefix of every consistent
play has at least that weight. Dividing by its length and taking a
lower limit proves Max's assertion.

For the other band, every low coordinate is below the upper boundary.
At a low Min vertex fix an edge attaining the minimum in $F_i(h)$,
and at a low Max vertex permit all outgoing edges. The supersolution
inequality gives
\begin{equation}
                         h_i\ge\widehat w(e)+h_j
 \label{eq:low-edge}
\end{equation}
for each permitted edge. Its head satisfies
$h_j\le(n+1)W<D_0-nW$, so it is low. These choices, extended
arbitrarily elsewhere, give a positional Min strategy preserving
the low band.

Every simple cycle of permitted edges now has nonpositive
transformed sum and hence original integer sum at most $-1$.
Deleting simple cycles from a prefix of length $T$ leaves at most
$n-1$ edges. Each deleted cycle has at most $n$ edges, so the
number removed is at least $(T-(n-1))/n$. The original prefix
weight is at most
\[
                       -\frac{T-(n-1)}n+(n-1)W_0.
\]
Its upper limiting average is at most $-1/n$. The prefix estimates
on both sides apply to every consistent edge sequence, independently
of how the opposing player uses the history.
\end{proof}

\subsection{Simultaneous boundary obligations}

The recursion works on inputs
\begin{equation}
 A\in\Z^V,\qquad D=2^d\ (d\in\Z_{\ge0}),\qquad
 a,b\in\Q_{>0}^V.
 \label{eq:recursive-input}
\end{equation}
Its box is $[A,u]$ with $u=A+D$, and its claim margin is
$P=D/32$. For $S\subseteq V$ let
$a(S)=\sum_{i\in S}a_i$ and $b(S)=\sum_{i\in S}b_i$.
The positive rational vectors $a,b$ are called \emph{masses}.

\begin{definition}[Witnesses and their claims]
\label{def:witnesses}
A \emph{Max witness} for $(A,D,a,b)$ is a subsolution
$x\in[A,u]$ with support $S_x=\{i:x_i>A_i\}$ and
$a(S_x)\le1$. It \emph{claims} each $i$ with $x_i\ge u_i-P$.
A \emph{Min witness} is a supersolution $y\in[A,u]$ with support
$S_y=\{i:y_i<u_i\}$ and $b(S_y)\le1$. It claims each $i$
with $y_i\le A_i+P$. A witness is \emph{relevant} when it claims
at least one index.
\end{definition}

Since $P<D$, every claim belongs to its witness's support. Claims
of opposite types cannot coincide: otherwise box comparison would
give $u_i-P\le x_i\le y_i\le A_i+P$, contradicting $2P<D$.
Masses limit the witnesses covered by an obligation; the graph itself
always retains every vertex and edge.

\begin{theorem}[Simultaneous labelling]
\label{thm:labelling}
For every input in Equation~\eqref{eq:recursive-input}, the
procedure $\Label(A,D,a,b)$ of Section~\ref{sec:algorithm}
terminates and returns a vector in $\{+,-\}^V$. It assigns $+$
to every index claimed by any Max witness and $-$ to every index
claimed by any Min witness.
\end{theorem}

All witnesses of a call are covered by this one output. They are
neither enumerated nor supplied to the procedure. On vertices with
no claim the theorem permits either label.

Apply this guarantee to
\begin{equation}
                  A=0,\qquad D=D_0,\qquad a_i=b_i=1/n.
 \label{eq:main-call}
\end{equation}
Every support then has mass at most one. The box solution $h$ is
therefore a witness of both types. Moreover
$P=D_0/32\ge nW$, so its Max claims include all high vertices and
its Min claims include all low vertices. By
Proposition~\ref{prop:game-reduction}, a labelling satisfying
Theorem~\ref{thm:labelling} returns the exact winning set by taking
its positively labelled vertices. This uses the existence and
boundary location of $h$, without executing $\Basic$ on the wide
initial box.
